Ambiguity note: “$n$ has at most $r$ prime divisors” can mean either
- $\omega(n)\le r$ (at most $r$ distinct prime divisors), or
- $\Omega(n)\le r$ (at most $r$ prime factors counted with multiplicity).
The phrase “prime divisors” most naturally means $\omega(n)$. I record statements for $\Omega$ when they remain valid.

For $r\in\mathbb{N}$ define the set
\[
S_r := \{m\in\mathbb{N}: \exists k\in\mathbb{N}_{\ge 0},\ \exists n\in\mathbb{N}_{\ge 1}\text{ with } m=2^k+n \text{ and } \omega(n)\le r\}.
\]
(One may analogously define $S_r^{(\Omega)}$ using $\Omega(n)\le r$.)

Let $d(S_r)$ denote the (natural) density if the limit exists:
\[
d(S_r)=\lim_{N\to\infty}\frac{|S_r\cap\{1,2,\dots,N\}|}{N},
\]
and otherwise interpret “density at least $1-\epsilon$” as lower density $\underline d(S_r)=\liminf_{N\to\infty} |S_r\cap[1,N]|/N$.

Question: For every $\epsilon>0$, does there exist an integer $r=r(\epsilon)$ (written $r\ll_\epsilon 1$) such that $\underline d(S_r)\ge 1-\epsilon$?
